% ===========================================================
%  hea-bench: An open benchmark suite and reference baselines
%  for high-entropy alloy phase prediction
%
%  Drafted for SoftwareX (Elsevier). Build with pdfLaTeX:
%      pdflatex main && bibtex main && pdflatex main && pdflatex main
% ===========================================================

\documentclass[5p,times]{elsarticle}

% --- Encoding / fonts: pdfLaTeX-compatible only ---
\usepackage[T1]{fontenc}
\usepackage[utf8]{inputenc}
\usepackage{lmodern}
\usepackage{textcomp}     % for special characters: \texttimes, \textmu, etc.
\usepackage{textgreek}    % \textDelta, \textOmega, etc. without math-mode

% --- Math / units ---
\usepackage{amsmath, amssymb}
\usepackage{siunitx}
\sisetup{detect-all, range-phrase = --, range-units = single}

% --- Tables / figures ---
\usepackage{booktabs}
\usepackage{graphicx}
\usepackage{tabularx}
\usepackage{multirow}

% --- Misc ---
\usepackage[colorlinks=true, allcolors=blue]{hyperref}
\usepackage{url}
\usepackage{microtype}

% Convenient macros for compositions and descriptors.
% Use \ensuremath so they work in both text and math contexts.
\newcommand{\Smix}{\ensuremath{\Delta S_{\mathrm{mix}}}}
\newcommand{\Hmix}{\ensuremath{\Delta H_{\mathrm{mix}}}}
\newcommand{\Tm}{\ensuremath{T_{\mathrm{m}}}}
\newcommand{\Omegaparam}{\ensuremath{\Omega}}
\newcommand{\hea}{\texttt{hea-bench}\xspace}

\usepackage{xspace}

\journal{SoftwareX}

\begin{document}

\begin{frontmatter}

\title{%
  hea-bench: An open benchmark suite and reference baselines for high-entropy
  alloy phase prediction
}

\author[1]{David~Fieser\corref{cor1}}
\ead{dfieser9@gmail.com}
\cortext[cor1]{Corresponding author.}

\address[1]{Independent researcher; ORCID \href{https://orcid.org/0009-0007-5754-4331}{0009-0007-5754-4331}}

\begin{abstract}
Empirical phase-prediction rules for high-entropy alloys (HEAs) ---
the entropy criterion of Yeh, the atomic-size mismatch criterion of
Zhang, the valence-electron-concentration criterion of Guo and Liu,
and the \Omegaparam{}-parameter criterion of Yang and Zhang --- are
widely used as fast pre-screens for candidate compositions. Their
quoted accuracies, however, come from heterogeneous, mutually
incompatible datasets, and no published artifact treats them as
proper diagnostic classifiers on a common open benchmark. We present
\hea{}, a Python library and reproducible benchmark suite that
(i)~consolidates three primary open HEA-phase datasets (Borg
et~al. 2020; Pei et~al. 2020; Peivaste et~al. 2023) into a single
\num{7784}-composition resource with explicit per-row source
provenance and cross-source label-conflict tracking,
(ii)~implements all four canonical empirical rules as binary or
multi-class diagnostic classifiers with full diagnostic statistics
(accuracy, sensitivity, specificity, Wilson 95\% intervals,
receiver-operating-characteristic curves, Youden's $J$), and
(iii)~ships in two interchangeable user surfaces: a standard
\texttt{pip}-installable Python package and an in-browser frontend
that runs the same wheel via Pyodide / WebAssembly, requiring no
local install. On the consolidated benchmark, both binary rules
collapse to ``predict single-phase almost always'' (Youden's $J$ of
0.075 and 0.032 respectively), and a receiver-operating-characteristic
sweep of the Zhang $\delta$ threshold finds an optimum at
$\delta < 2.5\%$ rather than the canonical $\delta < 6.5\%$. We
release \hea{} 0.0.1 with 151 tests, including pinned regression
tests on every headline benchmark number, under the MIT License.
\end{abstract}

\begin{keyword}
high-entropy alloys \sep
multi-principal element alloys \sep
phase prediction \sep
benchmark suite \sep
empirical descriptors \sep
Miedema model \sep
Pyodide \sep
open data \sep
diagnostic classifiers
\end{keyword}

\end{frontmatter}

% ---------------------------------------------------------------------
% Required SoftwareX metadata block
% ---------------------------------------------------------------------

\section*{Required Metadata}

\begin{table}[h]
\caption{Code metadata.}
\centering
\begin{tabular}{p{0.39\linewidth} p{0.55\linewidth}}
\toprule
\textbf{Nr.} & \textbf{Code metadata description} \\
\midrule
C1 Current code version & 0.0.1 \\
C2 Permanent link to code/repository & to be assigned via Zenodo on first tagged release \\
C3 Permanent link to reproducible capsule & N/A (the package is dependency-free in its core; reproducible via \texttt{pip install hea-bench}) \\
C4 Legal code license & MIT \\
C5 Code versioning system used & git \\
C6 Software code languages, tools, and services used & Python ($\geq$3.10), Pyodide (browser deployment) \\
C7 Compilation requirements, operating environments \& dependencies & Pure Python, dependency-free core; optional extras: \texttt{pandas}, \texttt{matplotlib}, \texttt{pytest} \\
C8 If available, link to developer documentation/manual & \texttt{README.md} in the repository \\
C9 Support email for questions & \texttt{dfieser9@gmail.com} \\
\bottomrule
\end{tabular}
\end{table}

% =====================================================================
\section{Motivation and significance}
% =====================================================================

High-entropy alloys (HEAs), introduced independently by Cantor
et~al.~\cite{cantor2004} and Yeh et~al.~\cite{yeh2004} in 2004, are
multi-component alloys (typically five or more principal elements
near equiatomic composition) in which the configurational mixing
entropy is large enough to stabilize disordered solid solutions
over the intermetallic compounds predicted by classical
Hume-Rothery rules. The compositional space is vast --- combinatorially
larger than that of conventional alloys --- which makes
\emph{pre-screening} a practical necessity before any expensive
density-functional-theory, CALPHAD, or experimental campaign.

Four empirical phase-prediction rules dominate the literature:
the configurational mixing entropy criterion of Yeh
($\Smix > 1.5 R$~\cite{yeh2004}); the atomic-size mismatch criterion
of Zhang ($\delta < 6.5\%$~\cite{zhang2008}); the valence-electron
concentration criterion of Guo and Liu (FCC favored at $\mathrm{VEC}
\geq 8.0$, BCC at $\mathrm{VEC} < 6.87$~\cite{guo2011}); and the
$\Omegaparam$-parameter criterion of Yang and Zhang
($\Omegaparam = \Tm \cdot \Smix / |\Hmix| > 1.1$~\cite{yang2012}).
Each is supported in its original publication by classification
rates of roughly $80\%$ on the dataset available at the time.

Three gaps in the open literature motivate this work.
\textbf{First,} every published evaluation of these rules uses a
different dataset, with different preprocessing, dedup, and phase
labeling. Head-to-head comparison is therefore impossible. To our
knowledge there is no published, open, MatBench-style benchmark
suite~\cite{dunn2020matbench} for HEA phase classification.
\textbf{Second,} the rules are typically reported as ``accuracies''
of the form $X / N$, without the diagnostic-classifier machinery
(sensitivity, specificity, Wilson intervals, ROC analysis,
Youden's~$J$) that would expose where they actually generalize.
\textbf{Third,} the only field-standard interactive tool for
Miedema-style enthalpy calculations is the Windows-only,
registration-gated MiedCalc~\cite{zhang2016miedema}; no
zero-install open analogue exists.

\hea{} addresses all three. We consolidate three primary open
HEA-phase datasets into a single $\num{7784}$-composition resource
with full provenance, implement the four canonical rules as proper
diagnostic classifiers, and ship the result in two equivalent
surfaces: a standard Python library and an in-browser frontend that
runs the same wheel via Pyodide.

% =====================================================================
\section{Software description}
% =====================================================================

\subsection{Software architecture}

\hea{} is organized into four orthogonal sub-packages plus a thin
command-line interface:

\begin{itemize}
  \item \texttt{hea\_bench.composition} --- formula parsing and
    composition normalization. A single regular-expression-based
    parser handles all three source formula conventions (Borg's
    space-separated proportional amounts, Pei's compact form,
    Peivaste's bare element strings).

  \item \texttt{hea\_bench.descriptors} --- the six empirical
    descriptors used by the rule set: $\Smix$, $\delta$, VEC,
    $\Tm$, $\Hmix$ (Miedema, multi-component~\cite{miedema1973,miedema1980}),
    and $\Omegaparam$~\cite{yang2012}. Per-element parameters live
    in \texttt{descriptors/data/elemental.py} (24 elements);
    binary mixing-enthalpy pairs are loaded from a vendored copy of
    matminer's~\cite{ward2018} Miedema pair table (BSD-3-Clause,
    license preserved alongside the data file), which itself
    derives from the de Boer parameter set~\cite{deboer1988} and
    is consistent with Takeuchi and Inoue's open
    tabulation~\cite{takeuchi2005}.

  \item \texttt{hea\_bench.rules} --- the four canonical empirical
    rules wrapped as classifiers, each in its own module
    (\texttt{yeh\_smix}, \texttt{zhang\_delta}, \texttt{guo\_vec},
    \texttt{yang\_omega}) with a tunable threshold where the
    literature value is the default.

  \item \texttt{hea\_bench.classifiers.diagnostic\_stats} --- pure
    stdlib implementations of Wilson score intervals~\cite{wilson1927},
    confusion matrices, binary classifier statistics, and ROC
    sweeps with Youden's~$J$~\cite{youden1950}.

  \item \texttt{hea\_bench.benchmark} --- per-source loaders (Borg,
    Pei, Peivaste), the cross-source consolidator, and a coverage
    analysis that quantifies which fraction of the consolidated
    dataset is scorable by the current elemental data tables.

  \item \texttt{hea\_bench.evaluate} --- the orchestrator that runs
    each rule against the consolidated benchmark and writes
    \texttt{rule\_baselines.json} alongside the consolidated CSV.
\end{itemize}

Two end-user surfaces consume this library. The Python surface is
a standard \texttt{pip}-installable package
(\texttt{pip install hea-bench}). The browser surface is a single
HTML page that loads Pyodide~\cite{pyodide} and invokes
\texttt{micropip.install()} on the same wheel that
\texttt{pip} would deliver; descriptor and rule calls happen
client-side in the user's browser tab, with no server-side
component. \emph{The library has one canonical implementation;
both surfaces run it identically.}

\subsection{Software functionalities}

\begin{enumerate}
  \item Parse arbitrary composition formulas, normalize to mole
    fractions, and report the six standard descriptors.
  \item Apply each of the four canonical empirical rules and obtain
    a predicted phase classification.
  \item Score a user-supplied rule (or composition set) against the
    consolidated v0.1.0 benchmark with full diagnostic statistics.
  \item Sweep a rule threshold and report the receiver-operating-
    characteristic curve, optimal accuracy threshold, and optimal
    Youden's-$J$ threshold.
  \item Analyse element coverage of an alloy set against the
    descriptor data tables and report next-best-add elements.
  \item Re-consolidate the per-source data via the bundled loaders
    and dedupe / conflict-detection logic.
\end{enumerate}

\subsection{The consolidated benchmark}

\hea{} v0.1.0 ships a single canonical benchmark at
\texttt{data/consolidated/v0.1.0/consolidated.csv}
($\num{7784}$ unique compositions). It is built by running the
per-source loaders on three open datasets:

\begin{itemize}
  \item Borg et~al.~\cite{borg2020} ($\num{740}$ unique
    \texttt{(formula, processing)} alloys; CC-BY-4.0; mirrored
    verbatim with SHA-256 stamp);
  \item Pei et~al.~\cite{pei2020} ($\num{1209}$ unique
    formulas; CC-BY-4.0; mirrored);
  \item Peivaste et~al.~\cite{peivaste2023} ($\num{7747}$ unique
    formulas after dedup of an upstream-generated combinatorial
    sweep; the upstream data repository declares no license and is
    therefore pointer-only, fetched on user request).
\end{itemize}

Compositions are merged on a 4-decimal-precision composition
key. Borg's processing column is preserved as a side-channel rather
than included in the join key. When two sources assign different
canonical labels to the same composition, the row is flagged
($\texttt{has\_conflict} = 1$) and the per-source labels are
preserved verbatim. The v0.1.0 release contains \num{100} such
cross-source conflicts (\num{1.3}\% of the benchmark), which are
themselves a useful artifact: they enumerate compositions on
which the open HEA literature does not agree.

The canonical phase taxonomy is deliberately coarse for v0.1.0 ---
$\{\mathrm{BCC},\, \mathrm{FCC},\, \mathrm{HCP},\,
\mathrm{multi\text{-}phase}\}$ --- so that all three sources map
into it without lossy guessing. Finer per-source labels (Borg's
microstructure column with $30+$ sublabels, Peivaste's 12-category
column distinguishing intermetallics and amorphous alloys) are
preserved as side-channel columns.

% =====================================================================
\section{Illustrative examples}
% =====================================================================

\subsection{Cantor alloy, end to end}

The canonical equimolar CoCrFeMnNi alloy serves as the field's
reference HEA and as our regression sanity check. With \hea{}:

\begin{verbatim}
import hea_bench
cantor = hea_bench.parse_formula("CoCrFeMnNi")

hea_bench.smix(cantor)               # 13.381 J/(mol K)  = R ln 5
hea_bench.delta(cantor)              # 3.164  %
hea_bench.vec(cantor)                # 8.000
hea_bench.mixing_enthalpy(cantor)    # -4.16  kJ/mol
hea_bench.omega(cantor)              # 5.794
\end{verbatim}

All four canonical rules classify Cantor as HEA-class,
single-phase, FCC --- consistent with the original experimental
characterization~\cite{cantor2004}. Every value above is pinned in
the regression test suite.

\subsection{Reference rule baselines on the consolidated benchmark}

The orchestrator \texttt{python -m hea\_bench.evaluate} produces
the headline numbers in Table~\ref{tab:rule-baselines}. These are
the reference values our test suite pins; any future drift in
dataset, descriptor code, vendored pair-enthalpy table, or rule
thresholds also breaks the build.

\begin{table*}[t]
\centering
\caption{Reference baselines for the four canonical empirical HEA
phase-prediction rules, evaluated as diagnostic classifiers against
the \hea{}~v0.1.0 consolidated benchmark.
$n$ is the number of compositions in the benchmark scorable by the
rule's required descriptor stack (the smaller numbers for the
binary rules reflect compositions fully covered by both
$\mathrm{ELEMENTAL\_DATA}$ and the Miedema pair table; the Guo--Liu
rule is stratified to single-phase observations only, as discussed
in the text). Wilson 95\% intervals are reported for accuracy.
The two binary rules have Youden's $J$ close to zero, indicating
that they collapse to ``predict single-phase almost always'' on
this benchmark; the Guo--Liu rule shows a strong FCC bias.}
\label{tab:rule-baselines}
\begin{tabular}{lrrrrr}
\toprule
Rule & $n$ & Accuracy & Sensitivity & Specificity & Youden's $J$ \\
\midrule
Zhang $\delta < \SI{6.5}{\percent}$ & \num{6651} & \SI{56.7}{\percent} [\SI{55.5}{\percent}, \SI{57.9}{\percent}] & \SI{99.0}{\percent} & \SI{8.5}{\percent}  & 0.075 \\
Yang  $\Omegaparam > 1.1$           & \num{6651} & \SI{54.4}{\percent} [\SI{53.2}{\percent}, \SI{55.6}{\percent}] & \SI{95.8}{\percent} & \SI{7.4}{\percent}  & 0.032 \\
Guo--Liu VEC (stratified to FCC$|$BCC observed) & \num{3463} & \SI{66.9}{\percent} & \SI{92.3}{\percent} (FCC) / \SI{48.3}{\percent} (BCC) & --- & --- \\
\bottomrule
\end{tabular}
\end{table*}

The Yeh $\Smix$ rule is descriptive rather than predictive;
$47.0\%$ of the benchmark passes the $\Smix > 1.5 R$ HEA-class
threshold, $36.9\%$ falls in the MEA range ($R \leq \Smix
\leq 1.5R$), and $16.1\%$ is dilute.

\subsection{ROC recalibration of the Zhang $\delta$ rule}

Treating the rules as proper classifiers admits a question the
canonical literature does not address: \emph{what threshold would
maximize Youden's $J$ on this benchmark?} Sweeping the Zhang
$\delta$ threshold from \SI{1}{\percent} to \SI{12}{\percent} in
\SI{0.1}{\percent} steps yields the result in
Table~\ref{tab:roc-zhang}.

\begin{table}[h]
\centering
\caption{ROC sweep of the Zhang $\delta$ rule on the \hea{}~v0.1.0
benchmark. The canonical \SI{6.5}{\percent} threshold yields
Youden's $J = 0.075$; the dataset-optimal $J$ is \num{0.113} at
$\delta < \SI{2.5}{\percent}$, $\sim$2.6$\times$ tighter.}
\label{tab:roc-zhang}
\begin{tabular}{lrrr}
\toprule
Threshold      & Sens. & Spec. & Youden's $J$ \\
\midrule
\SI{6.5}{\percent} (canonical) & \SI{99.0}{\percent} & \SI{8.5}{\percent}  & 0.075 \\
\SI{6.6}{\percent} (max acc.)  & \SI{99.5}{\percent} & \SI{8.2}{\percent}  & 0.075 \\
\SI{2.5}{\percent} (max $J$)   & \SI{23.6}{\percent} & \SI{87.7}{\percent} & \textbf{0.113} \\
\bottomrule
\end{tabular}
\end{table}

The canonical \SI{6.5}{\percent} threshold is heavily skewed
toward sensitivity at the expense of specificity. The
balanced-optimal threshold on the consolidated benchmark is
$\sim$2.6$\times$ tighter. Whether $J = 0.113$ constitutes a
``useful'' classifier is a separate question --- it does not ---
but the recalibration finding itself is a direct, defensible
output of treating the rule as a diagnostic classifier rather than
a fixed cutoff.

% =====================================================================
\section{Impact}
% =====================================================================

\hea{} consolidates an open benchmark that previously did not exist
in unified form, and supplies the diagnostic-classifier machinery
that previously was not applied to the canonical HEA empirical
rules in the open literature. We expect three classes of impact.

\textbf{First, calibration of expectations.} The collapse of both
binary rules to near-zero Youden's $J$ on the consolidated
benchmark is a quantitative caveat for downstream users of those
rules. The widely-quoted ``80\% accuracy'' figures in the original
papers reflect different, smaller datasets and a different choice
of evaluation metric; on a consolidated open dataset the rules
behave very differently.

\textbf{Second, a substrate for future methods work.} Any
candidate ML model, CALPHAD-derived classifier, or new empirical
descriptor can now be scored against the same dataset and the
same diagnostic-stats vocabulary as the four canonical rules.
The result is a head-to-head comparison rather than incommensurable
single-paper accuracies.

\textbf{Third, accessibility.} The Pyodide front-end lets any
researcher or student type a composition into a browser tab and
see the same numbers the published library would produce on a
local install --- with the wheel itself loaded into the
browser, not a re-implementation. This removes the install-friction
barrier without sacrificing the library-as-canonical-implementation
guarantee.

% =====================================================================
\section{Conclusions}
% =====================================================================

We have presented \hea{} 0.0.1, an open Python library and
benchmark suite for high-entropy alloy phase prediction.
It ships a $\num{7784}$-composition consolidated benchmark drawn
from three primary open sources with explicit cross-source
provenance and conflict tracking; reference-baseline diagnostic-
classifier evaluations of the four canonical empirical rules; a
ROC-based recalibration finding for the Zhang $\delta$ rule
showing a dataset-optimal threshold of $\delta < \SI{2.5}{\percent}$
rather than the canonical $\SI{6.5}{\percent}$; and an in-browser
Pyodide front-end that runs the same Python wheel a local
\texttt{pip install} would deliver. The package is MIT-licensed,
covered by $151$ tests including pinned regression tests on every
headline benchmark number, and intended to serve as the substrate
on which subsequent ML and CALPHAD methods for HEA phase
prediction can be compared.

Planned future work includes expanding the elemental-property
table to lift benchmark coverage from \SI{86.7}{\percent} to
$\sim$95\% (the additions are low-effort because the corresponding
pair-enthalpy entries are already present in the vendored Miedema
table), and a v0.2.0 release that exposes the cross-source label
conflicts as a structured side-dataset for downstream
disagreement-resolution work.

% ---------------------------------------------------------------------
\section*{Conflict of interest}
The author declares no competing financial or non-financial
interests.

\section*{Acknowledgements}
We thank the matminer maintainers~\cite{ward2018} for the open
Miedema pair-enthalpy table vendored under BSD-3-Clause; the
Pyodide project~\cite{pyodide} for the browser-side Python runtime;
and the authors of the three primary source datasets~\cite{borg2020,
pei2020, peivaste2023} for releasing data under reusable terms.
Code scaffolding and documentation drafts were prepared with the
assistance of large language models (Anthropic Claude); all
numerical parameters, formulae, threshold values, and benchmark
results were independently verified against the cited primary
sources or computed from documented inputs by the author.

% ---------------------------------------------------------------------
\bibliographystyle{elsarticle-num}
\bibliography{refs}

\end{document}
